\title{Parameter-Efficient Orthogonal Finetuning via Factorization}

\begin{abstract}
Large foundation models are becoming increasingly prevalent, but training them from the ground up remains prohibitively costly. Therefore, adapting these potent models efficiently to downstream tasks is gaining significance. In this work, we explore a principled finetuning strategy — Orthogonal Finetuning (OFT) — for downstream adaptation. Although OFT exhibits strong generalization capabilities, it still requires a relatively large number of trainable parameters due to the high dimensionality associated with orthogonal matrices. To tackle this, we first analyze OFT through the lens of information transmission and then identify several essential criteria that enhance parameter efficiency. Drawing inspiration from the Cooley-Tukey fast Fourier transform algorithm, which facilitates efficient information transfer, we introduce a parameterization of orthogonal matrices using butterfly structures. Integrating this parameterization with OFT, we develop a novel parameter-efficient finetuning method named Orthogonal Butterfly (BOFT). BOFT generalizes OFT as a particular instance, thereby establishing a broader orthogonal finetuning framework. Finally, we perform comprehensive empirical evaluations by adapting large vision transformers, substantial language models, and text-to-image diffusion models across diverse vision and language downstream tasks.
\end{abstract}

\section{Introduction}

Recent models like ChatGPT~\cite{brown2020language,chen2021evaluating} and Stable Diffusion~\cite{rombach2022high} showcase the impressive generalization capabilities of large foundation models. The swift improvement in these models’ performance comes with a substantial growth in parameter count (e.g., GPT-3 contains approximately 175 billion parameters). Consequently, training such foundation models from scratch has become increasingly challenging. Progress in the field therefore hinges on the ability to adapt these models without full retraining. In other words, it is critical to efficiently tailor existing foundation models to new downstream tasks. There are three main approaches for efficient task adaptation: \emph{model finetuning}~\cite{hulora2022,zaken2022bitfit,guo2020parameter,raffel2020exploring,qiu2023controlling,zhang2022adaptive,chavan2023one}, where the original model architecture is maintained and only a subset of parameters is finetuned; \emph{adapter tuning}~\cite{rebuffi2017learning,houlsby2019parameter,pfeiffer2020adapterfusion,lin2020exploring,he2021towards}, which involves adding trainable modules to the existing model; and \emph{prompt tuning}~\cite{li2021prefix,lester2021power}, which prepends trainable tokens to the input. Among these techniques, model finetuning stands out as a straightforward yet effective method and, importantly, does not introduce additional inference latency.

The core idea behind model finetuning is to maintain similarity between the pretrained model and the finetuned model according to certain metrics, thereby preserving the knowledge acquired during pretraining. For example, many current finetuning techniques use a small learning rate as a heuristic to limit the divergence between the pretrained and finetuned models. Considering the structured characteristics of weight matrices, a more systematic strategy aims to retain the relational information within the weight matrices, specifically the pairwise angles between neurons. This concept gave rise to a new finetuning framework called Orthogonal Finetuning~(OFT)~\cite{qiu2023controlling}, where neurons within the same layer undergo transformation via a shared orthogonal matrix, ensuring that the pairwise angles between neurons remain theoretically intact throughout finetuning. Although OFT has shown promising results in terms of generalization and convergence when finetuning text-to-image diffusion models~\cite{qiu2023controlling}, the number of trainable parameters can be substantial due to the high dimensionality of orthogonal matrices. To mitigate this, OFT employs a block-diagonal structure to reduce the parameter count. Nonetheless, this gain in parameter efficiency comes with trade-offs — the orthogonal matrix adopts a fixed sparsity pattern, and the orthogonal transformation is applied independently across different blocks. This block-diagonal sparsity, while parameter-saving, may introduce unwanted inductive biases (e.g., it limits expressiveness by restricting the orthogonal matrix’s ability to approximate classical linear transformations), and importantly, the method for identifying an optimal sparsity pattern remains unclear.

The main challenge in tackling this issue is to produce a dense orthogonal matrix while maintaining parameter efficiency. Although this might initially appear impractical since a $d$-dimensional dense orthogonal matrix demands $\mathcal{O}(d^2)$ parameters, we adopt an innovative strategy by composing a dense orthogonal matrix from multiple factorized sparse matrices. This method is based on the insight that the number of trainable parameters can be minimized by exchanging computational time for memory space. Because we represent the orthogonal matrix as a product of sparse matrices, this multiplication must be performed repeatedly during each training iteration. To better understand the matrix factorization, we frame the generation of a dense orthogonal matrix as a problem of information transmission. Within this framework, producing a dense orthogonal matrix via multiplication of sparse matrices can be viewed as transmitting information across a grid-structured graph. This perspective allows us to devise numerous ways to execute sparse matrix factorization that constrain the number of trainable parameters while retaining sufficient expressiveness to generate dense matrices.

To enhance parameter efficiency in our information transmission model, we draw inspiration from the butterfly structures found in the Cooley-Tukey fast Fourier transform algorithm~\cite{cooley1965algorithm}, which facilitates efficient information exchange among different nodes~\cite{klasing1994broadcasting}. In particular, the butterfly graph inherent in the Cooley-Tukey algorithm naturally suggests a method for performing sparse matrix factorization, termed butterfly factorization. Utilizing butterfly factorization, we can produce a $d$-dimensional dense matrix as a product of $\mathcal{O}(\log d)$ sparse matrices, each containing $\mathcal{O}(d)$ non-zero elements. By guaranteeing orthogonality for each sparse matrix, we achieve an efficient orthogonal parameterization requiring only $\mathcal{O}(d\log d)$ parameters, which represents a substantial reduction compared to the original parameter count. Leveraging this compact orthogonal parameterization, we introduce a novel parameter-efficient fine-tuning technique called Orthogonal Butterfly (BOFT). BOFT generalizes OFT as a special instance and establishes a broad orthogonal fine-tuning framework. A common feature shared by both the block-diagonal structure and the butterfly structure is sparsity. Both introduce sparsity into orthogonal matrices to decrease the effective number of trainable parameters. It is insightful to compare our approach with the low-rank structure used in LoRA~\cite{hulora2022}; both low-rank and sparse matrices constitute major categories of structured matrices~\cite{candes2011robust} that enable parameter efficiency.

In contrast to the block-diagonal structure employed by OFT to balance expressivity and regularity, BOFT utilizes the butterfly structure to create a smoother interpolation between matrices from the full orthogonal group (expressivity) and identity matrices (regularity). This approach allows us to identify a more compact hypothesis space within the orthogonal group suitable for downstream tasks. Given the extensive use of the butterfly structure in many rapid linear transforms, such as the discrete Fourier and discrete cosine transforms, we contend that our structured method for parameter efficiency introduces an implicit inductive bias that enhances generalizability and mitigates overfitting. Our reasoning is based on the fact that the butterfly structure can readily reconstruct numerous classic linear transforms, a capability that the block-diagonal structure in OFT lacks. Our main contributions are summarized as follows:

\begin{itemize}[leftmargin=*,nosep]
    \item We explore the challenge of parameter efficiency in orthogonal finetuning through a novel information transmission framework. This framework recasts the problem of designing a parameter-efficient dense orthogonal matrix as one of information transmission within a grid-structured graph.
    \item Drawing inspiration from the butterfly structures found in the Cooley-Tukey algorithm, we introduce Orthogonal Butterfly, a parameter-efficient orthogonal finetuning technique, situated within the information transmission framework.
    \item We offer theoretical insights explaining why BOFT can drastically reduce the number of trainable parameters while maintaining strong expressivity and training stability. Owing to the matrix factorization, BOFT also exhibits an intriguing weight interpolation property (refer to Figure~\ref{fig:boft_interpolation}).
    \item For the first time, we extend orthogonal finetuning beyond controllable text-to-image generation~\cite{qiu2023controlling}, demonstrating its broad potential as a universal model finetuning method. To this end, we apply BOFT across various adaptation scenarios spanning computer vision and natural language processing. BOFT surpasses existing state-of-the-art approaches by a significant margin, confirming its superior parameter efficiency and generalization capacity.
\end{itemize}

\section{Related Work}

\textbf{Parameter-efficient finetuning (PEFT)}. As foundation models (\emph{e.g.}, \cite{brown2020language,radford2021learning,rombach2022high,kirillov2023segment}) continue to grow in size and capability, there has been significant interest in fine-tuning these models for downstream tasks in a parameter-efficient manner~\cite{treviso2023efficient,ding2023parameter,lialin2023scaling}. Among the various PEFT techniques~\cite{houlsby2019parameter,aghajanyan2020intrinsic,hulora2022,edalati2022krona,wang2022adamix,gheini2021cross,zaken2022bitfit,guo2020parameter,sung2021training,ansell2022composable,lester2021power,li2021prefix,vu2022spot,he2021towards,mao2021unipelt,karimi2021compacter,liu2022few,sung2022lst,chen2022parameter,jia2022visual,chen2022adaptformer,zhang2022neural,jie2023fact,lian2022scaling,luo2023towards}, reparameterization-based approaches~\cite{aghajanyan2020intrinsic,hulora2022,edalati2022krona} are most closely related to our study. LoRA~\cite{hulora2022} modifies the pretrained weight matrix by adding the product of two low-rank matrices, showing promising results on natural language tasks. Because the rank of all added matrices is fixed in LoRA, several methods \cite{zhang2022adaptive,valipour2022dylora,zhang2023increlora} dynamically vary the rank across layers to better allocate the parameter budget. Due to its straightforwardness, this low-rank weight reparameterization has become highly popular \cite{dettmers2023qlora,zi2023delta,chavan2023one}. Motivated by how hyperspherical energy relates to generalization~\cite{liu2018learning,liu2021orthogonal}, \cite{qiu2023controlling} introduced orthogonal finetuning (OFT), an alternative yet effective technique for fine-tuning text-to-image diffusion models. OFT specifically learns an orthogonal matrix to transform neurons within the same layer, achieving superior generalization and consistently more stable training compared to LoRA. Although OFT performs well, it generally involves more trainable parameters than LoRA. Consequently, improving the parameter efficiency of OFT is a valuable objective. Furthermore, whether OFT can be extended to a broader range of adaptation tasks (beyond controlling text-to-image diffusion models~\cite{qiu2023controlling}) remains unclear. BOFT enhances OFT’s parameter efficiency through butterfly factorization. This advancement enables us to demonstrate the effectiveness of orthogonal finetuning across general adaptation scenarios.

\textbf{Butterfly structures}. The radix-2 Cooley-Tukey algorithm recursively decomposes the $N$-point discrete Fourier transform into two $\frac{N}{2}$-point discrete Fourier transforms, producing a butterfly structure expressible as a product of multiple sparse matrices (known as a butterfly matrix). Butterfly matrices have previously been employed to parameterize orthogonal matrices to circumvent pivoting in Gaussian elimination and enhance efficiency~\cite{parker1995random}, to stabilize recurrent neural network training~\cite{jing2017tunable}, and in kernel approximation~\cite{munkhoeva2018quadrature}. Works such as \cite{dao2019learning,dao2019kaleidoscope} leverage butterfly parameterizations to learn fast linear transforms. Similarly, \cite{chen2022pixelated,dao2022monarch} utilize butterfly matrices for efficient neural network training. Butterfly structures have also proven valuable in fast matrix-vector multiplication~\cite{michielssen1996multilevel,o2010algorithm}, data-sparse matrix approximation~\cite{li2015butterfly}, and network communication~\cite{klasing1994broadcasting,li2003linear,rajasingh2016transmission}. Differing from prior work, our focus lies in exploiting butterfly structures to improve the parameter efficiency of OFT.

\section{An Information Transmission View on Orthogonal Finetuning}

We begin with some foundational concepts of OFT. To fine-tune the pretrained weight matrix, OFT re-expresses the new weight matrix as the product of a learnable orthogonal matrix and the fixed pretrained weight matrix. Unlike LoRA, which updates the weights by adding a low-rank matrix, OFT updates weights via multiplication by an orthogonal matrix. For parameter efficiency, LoRA employs a low-rank structure, whereas the original OFT adopts a block-diagonal orthogonal structure~\cite{qiu2023controlling} (smaller diagonal block sizes lead to greater parameter efficiency for OFT). A straightforward comparison is shown in Figure~\ref{fig:comp_lora}. The rationale behind applying an orthogonal transformation to fine-tune the weight matrix is to maintain the pairwise angles between neurons~\cite{liu2018learning,liu2021orthogonal,qiu2023controlling}, thereby largely preserving the semantic knowledge acquired during pretraining. Specifically, OFT optimizes an orthogonal matrix $\bm{R}\in\mathbb{R}^{d\times d}$ for a pretrained linear layer $\bm{W}^0\in\mathbb{R}^{d\times n}$ and modifies the forward computation from $\bm{z}=(\bm{W}^0)^\top\bm{x}$ to $\bm{z}=(\bm{R}\bm{W}^0)^\top\bm{x}$, where $\bm{x}\in\mathbb{R}^{d}$ and $\bm{z}\in\mathbb{R}^{n}$ denote the input and output vectors, respectively. To ensure zero initialization, $\bm{R}$ is initialized as the identity matrix. To preserve the orthogonality of $\bm{R}$ throughout fine-tuning, we follow \cite{qiu2023controlling,liu2021orthogonal} and use the Cayley parameterization, i.e., $\bm{R}=(\bm{I}+\bm{Q})(\bm{I}-\bm{Q})^{-1}$, where $\bm{Q}$ is skew-symmetric with $\bm{Q}=-\bm{Q}^\top$. For parameter efficiency, the block-diagonal structure is introduced by expressing the orthogonal matrix $\bm{R}$ as $\text{diag}(\bm{R}_1,\bm{R}_2,\cdots,\bm{R}_r)$, where each $\bm{R}_i\in\mathcal{R}^{b\times b}$, $\forall i$, is a small orthogonal matrix and $br=d$. The parameter efficiency obtained via the block-diagonal structure comes with a trade-off—it assumes that the dimensions of a neuron (i.e., a column vector of $\bm{W}^0$) are partitioned into $r$ groups, with different orthogonal matrices applied separately to each group. Although the block-diagonal approach is empirically effective, dividing neuron dimensions into $r$ groups based solely on their indices lacks theoretical justification, which motivates the desirability of dense orthogonal matrices. This naturally raises the question: \emph{Is it possible to construct a dense orthogonal matrix without compromising parameter efficiency?}

To tackle this issue, we suggest constructing a dense orthogonal matrix as a product of several sparse orthogonal matrices. To offer a cohesive and intuitive framework for analyzing the sparsity pattern in orthogonal matrix factorization, we reformulate the challenge of generating a dense orthogonal matrix in OFT as an information transmission problem. Concretely, producing a dense matrix $\bm{R}\in\mathbb{R}^{d\times d}$ as a product of $m$ square matrices, \thickmuskip=2mu \medmuskip=2mu $\bm{R}=\bm{B}_m\bm{B}_{m-1}\cdots\bm{B}_1$, can be interpreted as sending information through a grid composed of $d\times (m+1)$ nodes, as shown in Figure~\ref{fig:info_trans}. This information transmission perspective stems from the observation that a $d$-dimensional dense square matrix can be seen as encoding dense connectivity from one set of $d$ nodes to another set of $d$ nodes. For the matrix $\bm{R}$, if the element $R_{ij}$ equals zero, it signifies that information from the $j$-th node cannot reach the $i$-th node. Conversely, if $R_{ij}$ is non-zero, information transmission is possible. Hence, expressing the dense matrix $\bm{R}$ as a product of matrices $\bm{B}_m\bm{B}_{m-1}\cdots\bm{B}_1$ can also be regarded as executing sequential information exchanges governed by the graphs associated with each $\bm{B}_i, \forall i$. The flow of information follows from $\bm{B}_1$ initially to $\bm{B}_m$ ultimately. As a specific illustration in Figure~\ref{fig:info_trans}, we examine the factorization $\bm{R}=\bm{B}_5\bm{B}_4\bm{B}_3\bm{B}_2\bm{B}_1$, visualizing its sparsity patterns and the corresponding induced graph. The graph depicted in Figure~\ref{fig:info_trans} results from unrolling the matrix multiplication. Within this induced graph, the matrix $\bm{B}_i$ functions as the connectivity matrix from the nodes at level $i$ to those at level $(i+1)$. More precisely, the $(j_1,j_2)$-th element of $\bm{B}_i$ indicates the presence of a directed edge from the $j_2$-th node at level $i$ to the $j_1$-th node at level $(i+1)$, where a zero value signifies absence of an edge. For the product $\bm{B}_5\bm{B}_4\bm{B}_3\bm{B}_2\bm{B}_1$ to form a dense matrix, it is necessary that every node at the first level can transmit information to all nodes at the sixth level. Restricting attention to $\bm{R}=\bm{B}_4\bm{B}_3\bm{B}_2\bm{B}_1$, which corresponds to source nodes at level one and receiver nodes at level five, reveals that information from node $1$ cannot reach node $3$. Consequently, the sparsity pattern of $\bm{B}_4\bm{B}_3\bm{B}_2\bm{B}_1$ contains a zero at position $(3,1)$. Similar relationships between the induced graph and sparsity pattern hold for $\bm{R}=\bm{B}_3\bm{B}_2\bm{B}_1$ and $\bm{R}=\bm{B}_2\bm{B}_1$. In general, for a matrix $\bm{R}\in\mathbb{R}^{d\times d}$ to be dense, the factorization matrices $\bm{B}_m,\cdots,\bm{B}_1$ must correspond to a set of directed edges arranged on a $d\times (m+1)$ grid, where each directed edge connects nodes only between consecutive levels (i.e., adjacent columns), ensuring that information from every node at the first level can propagate to every node at the $(m+1)$-th level.

The perspective of information transmission aids in better comprehending the sparsity pattern of factorization matrices in OFT. Figure~\ref{fig:blk_diag} illustrates the block-diagonal configuration of $\bm{R}$ in the original OFT. Although this block-diagonal arrangement reduces the number of trainable parameters, it does not yield a dense matrix $\bm{R}$. Our objective is to form a dense orthogonal matrix by combining $m$ sparse orthogonal matrices while using the fewest effective trainable parameters possible. From the information transmission viewpoint, the primary criteria for achieving this are \emph{(i)} \textbf{dense connectivity}: each node at the first level must have at least one path to every node at the last level, and \emph{(ii)} \textbf{minimum free edges}: the total count of edges should be minimized under the orthogonality constraint. It is important to note that orthogonality imposes a subtle restriction on the edges connecting adjacent levels. For instance, to ensure each matrix $\bm{B}_i$ is full-rank (a prerequisite for orthogonality), there must be $d$ edges creating a bijection between all nodes in the $i$-th level and those in the $(i+1)$-th level. This means the number of edges between adjacent levels is at least $d$ (\emph{e.g.}, 4 in the example shown in Figure~\ref{fig:info_trans}). These $d$ edges are essential for maintaining orthogonality and should not be included in the count of edges, as these elements are fixed and non-trainable (for example, in a $d \times d$ orthogonal matrix with $d$ non-zero elements, these entries can only be $\pm 1$). Since orthogonal matrices require fewer parameters than full matrices, the orthogonality condition introduces additional dependencies among edges. For example, in a $2 \times 2$ orthogonal matrix, a zero entry at position $(1,1)$ necessitates another zero at $(2,2)$ (\emph{i.e.}, the absence of one edge implies the absence of another). Thus, for each valid set of edge connections, orthogonality may sometimes enforce the addition or removal of edges. By representing the non-zero structure of the composed orthogonal matrix, the information transmission framework proves particularly beneficial for OFT, as we focus solely on the non-zero trainable components of $\bm{R}$, irrespective of their precise values.

A straightforward dense connection linking two levels requires $\mathcal{O}(d^2)$ edges (i.e., a single dense orthogonal matrix), resulting in $d^2 - d$ edges (for $d=4$, this equates to 12 edges). Figure~\ref{fig:info_trans} illustrates an example of a viable matrix factorization that uses a total of $10$ edges, which is actually fewer than what a single dense orthogonal matrix demands. This framework allows us to analyze the parameter efficiency of OFT from a graph-theoretic viewpoint, making it easy to identify feasible factorizations. We draw inspiration from an intriguing topology found in the Cooley-Tukey algorithm, known as butterfly graphs~\cite{cooley1965algorithm}, which efficiently connect $d$ source nodes to $d$ receiver nodes with only $\mathcal{O}(d \log d)$ edges. For instance, the topology depicted in Figure~\ref{fig:info_trans} uses $10$ edges to establish dense connectivity, whereas a butterfly network requires just $8$ edges. In the following, we explain how the butterfly structure can enhance parameter efficiency.

\section{Orthogonal Parameterization by Butterfly Factorization}

The butterfly structure was initially introduced in the Cooley-Tukey algorithm to perform fast Fourier transforms. In Fourier transforms, a localized change in the frequency domain induces a global effect in the spatial domain, which conceptually parallels our information transmission problem—where each node at the first level can send information to all nodes at the last level. The butterfly structure has also become a widely used computer network topology~\cite{solihin2015fundamentals,li2011linear} for effective information exchange. Assuming $k \geq 2$ is a power of two, we define the \emph{butterfly factor} $\bm{B}^F(k)$ as

\begin{equation}\label{eq:but_factor}
\begin{aligned}
    \bm{B}^F(k) = \begin{bmatrix}
    \text{diag}(\bm{d}_1) & \text{diag}(\bm{d}_2) \\
    \text{diag}(\bm{d}_3) & \text{diag}(\bm{d}_4)
    \end{bmatrix} \in \mathbb{R}^{k \times k},
\end{aligned}
\end{equation}

where $\bm{d}_i \in \mathbb{R}^{\frac{k}{2}}, \forall i$ are certain vectors. For $d = 2^N$, we then recursively define the $d$-dimensional \emph{butterfly matrix} $\bm{B}(d) \in \mathbb{R}^{d \times d}$ as

\begin{equation}
\begin{aligned}
    \bm{B}(d) = \tilde{\bm{B}}(d,d) \cdot \begin{bmatrix}
    \bm{B}_1(\frac{d}{2}) & \bm{0} \\
    \bm{0} & \bm{B}_2(\frac{d}{2})
    \end{bmatrix} = \tilde{\bm{B}}(d,d) \tilde{\bm{B}}(d,\frac{d}{2}) \cdots \tilde{\bm{B}}(d,2),
\end{aligned}
\end{equation}

where $\bm{B}_1(\frac{d}{2})$ and $\bm{B}_2(\frac{d}{2})$ represent two butterfly matrices of dimension $\frac{d}{2}$. We then introduce the \emph{butterfly component} as $\thickmuskip=2mu \medmuskip=2mu \tilde{\bm{B}}(d,k)=\text{diag}(\bm{B}^F_1(k),\cdots,\bm{B}^F_{d/k}(k))$, which is a block-diagonal matrix of dimension $d \times d$ with blocks of size $k$, where each $\bm{B}^F_i(k), \forall i$ denotes butterfly factors defined in Equation~\ref{eq:but_factor}. At this point, we can use the butterfly matrix to represent an orthogonal matrix. To do so, it suffices to ensure that every multiplicative factor $\tilde{\bm{B}}(d,k), \forall k$ within the butterfly matrix $\bm{B}(d)$ is orthogonal. We first consider the block-diagonal matrix $\tilde{\bm{B}}(d,2)$ having block size 2, and orthogonality of $\tilde{\bm{B}}(d,2)$ can be straightforwardly ensured by applying the Cayley transform (or a 2-dimensional rotation) to parameterize each block, as done in \cite{liu2021orthogonal,qiu2023controlling}. The sparsity pattern of each butterfly component corresponds to a permutation of the sparsity pattern of $\tilde{\bm{B}}(d,2)$, enabling all butterfly components to be easily parameterized as orthogonal matrices. This yields an efficient orthogonal matrix parameterization constructed from numerous $2 \times 2$ orthogonal matrices. Following \cite{chen2022pixelated}, we extend the butterfly matrices and define a block butterfly component $\tilde{\bm{B}}^b(d,k)$ where every entry in $\bm{d}_i, \forall i$ is replaced by a $b \times b$ matrix. To ensure that the block butterfly component $\tilde{\bm{B}}^b(d,2)$ is orthogonal, we parameterize each $2b \times 2b$ block matrix to be orthogonal. The non-zero pattern of other butterfly components $\tilde{\bm{B}}^b(d,k)$ with $k > 2$ is a block-wise permutation of the pattern in $\tilde{\bm{B}}^b(d,2)$ and hence can be similarly parameterized as orthogonal matrices. Putting these parts together, the forward pass in BOFT is

\begin{equation}
    \begin{aligned}\nonumber
    \bm{z}=\big(\bm{R}(m,b)\cdot\bm{W}^0\big)^\top\bm{x},~~~~\text{s.t.}~\bigg{\{}\bm{R}(m,b)=\prod_{i=1}^m \tilde{\bm{B}}^b_{(d,i)}~~\&~~\underbrace{\big(\tilde{\bm{B}}^b_{(d,j)}\big)^\top\tilde{\bm{B}}^b_{(d,j)}=\tilde{\bm{B}}^b_{(d,j)}\big(\tilde{\bm{B}}^b_{(d,j)}\big)^\top\!=\bm{I}_d}_{\forall j\in [1, m]}\bigg{\}},
    \end{aligned}
\end{equation}

Here, we denote $\tilde{\bm{B}}^b(d,2^{m - i+1})$ as $\tilde{\bm{B}}^b_{(d,i)}$ for simplicity, and $\bm{I}_d$ represents the identity matrix of dimension $d$. The orthogonal matrix $\bm{R}(m,b)\in\mathbb{R}^{d\times d}$ is constructed by multiplying several orthogonal butterfly components. For ease of notation, we refer to BOFT with $\bm{R}(m,\frac{b}{2})$ as BOFT($m$,$b$), where $b\geq 2$. When $m=1$, BOFT($1$,$b$) simplifies to the block-diagonal OFT~\cite{qiu2023controlling} having block size $b$. BOFT($1$,$d$) corresponds to the original OFT~\cite{qiu2023controlling} with an unconstrained full orthogonal matrix. BOFT($\log \frac{2d}{b}$,$b$) employs the block butterfly matrix $\bm{B}^{\frac{b}{2}}(d)$ as $\bm{R}$, resulting in a dense orthogonal matrix $\bm{R}$. Generally, BOFT($m$,$b$) involves $\frac{1}{2}(b-1)dm$ effective trainable parameters when fine-tuning a linear layer of dimension $d \times n$. If we choose the butterfly matrix parameters, i.e., $m=\log d$ and $b=2$, BOFT requires $\mathcal{O}(d\log d)$ parameters. In contrast, the original OFT with a fully dense orthogonal matrix uses $\mathcal{O}(d^2)$ parameters, while the block-diagonal OFT with block count $r$ uses $\mathcal{O}(bd)$. Consequently, the original OFT must set block size $b=d$ to produce a dense orthogonal matrix, whereas BOFT can achieve this with any choice of $b$.

\textbf{Identity initialization for BOFT}. Fine-tuning methods commonly begin from the exact pretrained model so that the fine-tuned model remains close to the pretrained one. For instance, LoRA initializes its low-rank weights to zero. In BOFT, we initialize every butterfly component as the identity matrix (i.e., the skew-symmetric matrix in the Cayley transform is set to zero).

\textbf{Multiplicative Dropout for BOFT}. LoRA~\cite{hulora2022} also incorporates a Dropout layer for the low-rank weight updates to mitigate overfitting. While conventional Dropout~\cite{srivastava2014dropout} naturally applies to LoRA, it is incompatible with BOFT due to its multiplicative weight update scheme. To overcome this, we introduce a multiplicative Dropout strategy for BOFT. Since the orthogonal matrix $\bm{R}(m,b)$ consists of $m$ orthogonal butterfly components that can be readily rearranged into $2b \times 2b$ block-diagonal orthogonal matrices, this multiplicative Dropout randomly selects $p_1$ percent of butterfly components and $p_2$ percent of diagonal blocks within each component, substituting these blocks with identity matrices.

\section{Intriguing Insights and Discussions}

\textbf{Expressiveness of BOFT}. The butterfly architecture combined with permutations can exactly replicate numerous well-known fast linear transforms~\cite{dao2019learning,dao2019kaleidoscope} (such as the fast Fourier transform and Hadamard transform). However, the capability of our orthogonal butterfly matrix to approximate an arbitrary orthogonal matrix remains uncertain. We begin by performing a simulation to approximate a randomly generated dense orthogonal matrix~\cite{anderson1987generation} of dimension $512\times 512$. The outcomes shown in Figure~\ref{fig:para_eff} represent averages over 10 different random seeds. The y-axis indicates the approximation error, while the x-axis corresponds to the number of effective trainable parameters. Each curve with the same color corresponds to BOFT with a fixed block size, and the leftmost point represents the error of BOFT($1$,$b$) (i.e., the standard block-diagonal OFT with block size $b$). Generally, BOFT offers superior parameter efficiency compared to OFT. For instance, BOFT($9$,$2$) exhibits greater expressiveness than BOFT($1$,$16$) while requiring significantly fewer parameters. Typically, BOFT configurations with smaller $b$ and larger $m$ demonstrate enhanced parameter efficiency; for example, BOFT($6$,$4$) uses considerably fewer parameters but achieves a comparable approximation error to BOFT($2$,$16$). Overall, the butterfly matrix represents a more structured subset of the orthogonal group (relative to the block-diagonal structure), which theoretically makes BOFT more expressive than OFT at the same block size.

\begin{theorem}[Expressiveness of BOFT]\label{thm:exp_boft}
    BOFT possesses greater expressiveness than OFT when using the same block size. To approximate every orthogonal matrix of dimension $d$ via butterfly matrices, one can multiply butterfly matrices as $\bm{B}_{d-1,1}(d)\bm{B}^\top_{d-1,2}(d)\cdots\bm{B}_{1,1}(d)\bm{B}^\top_{1,2}(d)$, where each $\bm{B}_{i,j}(d)$ is a butterfly matrix for all $i,j$.
\end{theorem}

Theorem~\ref{thm:exp_boft} motivates a straightforward extension of BOFT—the final orthogonal matrix generalizes to $\thickmuskip=2mu \medmuskip=2mu \bm{R}^G(m_1,b_1,m_2,b_2,l)=\bm{R}_{l,1}(m_1,b_1)\bm{R}_{l,2}^T(m_2,b_2)\cdots\bm{R}_{1,1}(m_1,b_1)\bm{R}_{1,2}^T(m_2,b_2)$, where $\bm{R}_{i,j}^T(m,b)$ denotes the orthogonal matrix utilized in BOFT. When $m_1=m_2=\log d$, $b_1=b_2=2$, and $l=d-1$, the matrix $\bm{R}^G(m_1,b_1,m_2,b_2,l)$ spans the entire orthogonal group. This composition of matrices is also known as the kaleidoscope hierarchy~\cite{dao2019kaleidoscope}. Nevertheless, it is important to note that increased expressiveness does not necessarily translate into improved finetuning performance, as full finetuning, despite its universal expressiveness, frequently results in suboptimal outcomes. Balancing expressiveness and regularization is crucial for the generalizability of model finetuning. BOFT expands the finetuning parameter space while incorporating structural priors, allowing us to better balance expressiveness with regularization.

\textbf{Spectral properties}. Orthogonal finetuning typically results in superior spectral characteristics compared to LoRA, because it exactly maintains the spectral norm of the pretrained weight matrix $\bm{W}^0$. This can be demonstrated through singular value decomposition: $\bm{W}^0=\bm{U}\bm{\Sigma}\bm{V}^\top$, where $\bm{U}$ and $\bm{V}$ are orthogonal matrices and $\bm{\Sigma}$ is a diagonal matrix containing singular values. Both OFT and BOFT apply an orthogonal matrix $\bm{R}$ on the left, producing the finetuned weights $\bm{R}\bm{U}\bm{\Sigma}\bm{V}^\top$, which leaves the largest singular value unchanged (i.e., the spectral norm of $\bm{W}^0$ remains intact). This preservation has been demonstrated to significantly enhance training stability and generalization~\cite{miyato2018spectral,yoshida2017spectral}. Additional intriguing mathematical characteristics are provided in Appendix~\ref{app:sn_property}.

\textbf{Orthogonal finetuning as learning bilinear similarity}. BOFT can be expressed as learning the bilinear similarity $\bm{w}^0_i\bm{R}\bm{x}$, where $\bm{w}^0_i$ denotes the $i$-th neuron (i.e., the $i$-th column vector) of the weight matrix $\bm{W}^0$. In this view, BOFT corresponds to learning a bilinear similarity matrix $\bm{R}$ subject to a strong constraint (namely, that $\bm{R}$ is orthogonal), which inherently relates to distance metric learning~\cite{xing2002distance} and bilinear forms~\cite{roman2005advanced}.

\textbf{Inductive bias and generalization in BOFT}. Because $\bm{R}(m,b)$ in BOFT typically constitutes a structured subset of the orthogonal group that restricts the hypothesis space, BOFT naturally introduces an inductive bias. We contend that the structured inductive bias derived from butterfly factorization enhances generalization, as it embodies the common structured patterns of many classic linear transforms~\cite{dao2019kaleidoscope}, including the discrete Fourier transform, discrete sine/cosine transform, and Hadamard transform. Furthermore, the sparse matrix factorization within BOFT may also contribute an implicit inductive bias~\cite{gunasekar2017implicit,li2018algorithmic,liu2019neural}.

\textbf{Comparison to butterfly-based sparse training}. Several studies~\cite{dao2019learning,dao2019kaleidoscope,chen2022pixelated,dao2022monarch} have explored sparse training using butterfly parameterization. These works generally concentrate on directly reparameterizing weight matrices with the butterfly parameterization and training neural networks from scratch. The approach in \cite{dao2022monarch} involves fine-tuning pretrained weights by initially projecting them onto a variant of butterfly matrices and subsequently optimizing these projected components for downstream tasks. In contrast, BOFT introduces a fundamentally different fine-tuning method that modifies the weights using layer-shared weight matrices.

\input{rebuttal-results}

\section{Concluding Remarks and Limitations}

In this paper, we introduce Orthogonal Butterfly, a versatile and parameter-efficient fine-tuning technique for foundation models that leverages the butterfly structure. The central idea behind enhancing parameter efficiency is to represent a dense orthogonal matrix as a product of multiple sparse orthogonal matrices. To facilitate identifying feasible matrix factorizations, we develop a graphical information transmission framework. Within this framework, we discover that the butterfly structure effectively satisfies our criteria for sparse orthogonal matrix factorization. We validate the practical performance of BOFT through fine-tuning experiments on large language models, large vision models, and text-to-image generative models. Our results further establish BOFT’s superiority as a general fine-tuning approach.

Although BOFT demonstrates strong empirical results, it is not without limitations. Because the final orthogonal matrix in BOFT is composed as a product of several orthogonal matrices, the training runtime overhead is somewhat greater than that of OFT. Enhancing BOFT’s training speed remains an open challenge. Fortunately, once fine-tuning is complete, the orthogonal matrices learned by BOFT can be directly merged into the pretrained model, incurring no extra latency during inference. Additionally, it is still uncertain whether the butterfly network represents the most efficient means of information transmission. Our information transmission framework opens avenues to draw insights from a different research field — computer networking — where the efficiency of network topologies for information flow is extensively studied. We anticipate that more efficient network architectures may be employed to construct dense orthogonal matrices in the future.

\end{document}